\subsection{Source Repository}
The Iteration 4 source, report, evidence and native Spark notebooks are available in the public GitHub repository:
\begin{center}
\url{https://github.com/QuyDuong106/INFOSYS722_BDAS}
\end{center}
The repository is organised around CRISP-DM. Its \path{iteration4_BDAS} folder contains the report, original data and retained Python prototype; \path{PySpark} contains the native Spark implementation, saved executions, visualization notebook and new feature-selection extension.

\subsection{Shared Google Drive Project}
The shareable project folder contains the Colab working files, input data and saved execution outputs:
\begin{center}
\href{https://drive.google.com/drive/folders/1ob80bZnuNHHwhGT3gY-HXrj1-4I0VZ_F?usp=sharing}{Open the shared Iteration 4 Google Drive folder}
\end{center}
In Colab, the project is mounted at \path{/content/drive/MyDrive/INFOSYS_722_Notebook/iteration4_BDAS}. The completed native holdout output is stored in its \path{results/spark_holdout} folder. The research LOPO workflow has a separate output folder and is not included as holdout evidence.

\subsection{Colab Entry Points}
\begin{itemize}
\item \href{https://colab.research.google.com/github/QuyDuong106/INFOSYS722_BDAS/blob/main/PySpark/run_spark_iterations_colab.ipynb}{Native Spark CART20/RFE20/RFE10 extension}: the primary workflow for new controlled Spark preparation experiments. The local \path{iteration4_BDAS/run_workflow.ipynb} invokes this same extension.
\item \href{https://colab.research.google.com/github/QuyDuong106/INFOSYS722_BDAS/blob/main/PySpark/run_spark_workflow.ipynb}{Saved Spark holdout execution}: retains the author's measured all-field logistic, random-forest and GBT results.
\item \href{https://colab.research.google.com/github/QuyDuong106/INFOSYS722_BDAS/blob/main/PySpark/spark_visualization_colab.ipynb}{Saved Spark visualization}: supports CRISP-DM data understanding through Spark-computed summaries and plots.
\item \path{iteration4_BDAS/run_python_workflow.ipynb}: retains the five-family Python/Optuna prototype separately from the primary Spark submission workflow.
\end{itemize}

\subsection{Copied Non-Spark Optuna LOPO Method}
The \href{https://github.com/QuyDuong106/INFOSYS722_BDAS/tree/main/iteration4_BDAS/LOPO}{LOPO source and completed evidence folder} contains the Research Paper nested Optuna implementation and its five completed single-seed evaluations. This is a non-Spark Python method, distinct from the Spark/CUDA notebook. Its README documents execution, preprocessing, stacking search scope and output provenance.

\subsection{Evidence Provenance}
The native Spark result tables and average confusion matrices use the latest completed raw \path{results/spark_holdout/metrics.json} output in the shared Drive folder, saved on 12 October 2026 and retrieved on 13 October 2026. The local snapshot in \path{results/spark_holdout_drive} retains the raw metrics, source notebook, source script and provenance metadata. The notebook invocation confirms training-only tuning, and validation confirms three models across all five seeds with 96 test predictions per evaluation. These fifteen evaluations use overlapping holdouts rather than independent cohorts. The earlier six-decimal notebook display remains archived in \path{results/spark_colab_reference}; it agrees with the raw output at the displayed precision. New Spark extension outputs are included only after the corresponding native execution is complete.

The visualization provenance record is retained in \path{results/spark_visualization_reference}. Outcome, categorical and other earlier figures were extracted from the saved GitHub Colab notebook at commit \texttt{99f8af2} (12 October 2026). The continuous histograms and outcome-grouped continuous distributions have been refreshed from the author's latest saved Colab execution on 15 October 2026, covering all thirty-one continuous predictors. Its exact notebook snapshot is \path{all_continuous_colab_2026-10-15.ipynb}, with notebook/figure hashes recorded in the provenance JSON. Spark computed these summaries; no local Spark rerun is claimed.
